\documentclass[10pt,letterpaper]{article} % Tamaño de fuente a 10pt
\usepackage[utf8]{inputenc}
\usepackage[spanish, es-tabla]{babel}
\usepackage[left=2.3cm, right=2.3cm, top=1.5cm, bottom=1.5cm]{geometry} 
\usepackage{setspace}
\usepackage{ragged2e}
\usepackage{enumitem}

% --- CONFIGURACIÓN PARA ELIMINAR SANGRÍAS ---
\setlength{\parindent}{0pt}
\setlength{\parskip}{0.8em} 
\setlist[itemize]{leftmargin=*, labelsep=0.5em, nosep} 
\makeatletter
\def\@flushglue{0pt plus 1fil}
\makeatother
% ----------------------------------------------------

\begin{document}

\singlespacing
\RaggedRight 

Bogotá D. C., 28 de marzo de 2026

\vspace{0.2cm}

Señores:\\
\textbf{CONSEJO CURRICULAR}\\
Tecnología en Electrónica Industrial\\
Ingeniería en Control y Automatización e Ingeniería en Telecomunicaciones\\
Universidad Distrital Francisco José de Caldas

\vspace{0.3cm}

\textbf{Asunto:} Solicitud de asignación de Jurados 

\vspace{0.3cm}

Respetados Consejeros:

Cordial saludo. La presente tiene como fin solicitar a ustedes la asignación de jurados para la sustentación del trabajo de grado titulado \textit{“Diseño e Implementación de una Plataforma de Estabilización para una Bicicleta Utilizando un Volante de Inercia”} con código de radicado \textbf{20252583027}, presentado por los estudiantes \textbf{Jeferson Hernán Hernández Moreno} con código \textbf{20201383013} y \textbf{Fabián Andrés Reyes Romero} con código \textbf{20201383011}, bajo la dirección del docente \textbf{Alexander Jiménez Triana, PhD}.

El trabajo de grado fue avalado por el Consejo Curricular conforme al Acuerdo No. 012 de diciembre 13 de 2022 bajo la modalidad de \textbf{Pasantía}.

\textbf{Objetivo General}

Diseñar una plataforma mecánica capaz de estabilizar una bicicleta durante un corto periodo de tiempo, simulando condiciones de equilibrio en estado de reposo o baja velocidad. La estabilización se logrará mediante un volante de inercia que, por medio de un sistema de control, aprovechará el momento angular generado por su rotación, junto con el efecto giroscópico, para oponerse activamente a las perturbaciones que puedan afectar el balance del sistema. Este diseño permitirá analizar y validar el comportamiento dinámico de la bicicleta frente a pequeñas variaciones.

\textbf{Objetivos Específicos}

\begin{itemize}
    \item Diseñar y construir un modelo físico funcional de la planta, que permita la integración de un sistema de control para analizar el comportamiento dinámico de una bicicleta en condiciones de equilibrio inestables.
    \item Diseñar e implementar un volante de inercia adecuado al sistema, realizando las modificaciones necesarias a la estructura de la bicicleta para mantener el centro de masa lo más centrado posible, optimizando así el control y la estabilidad.
    \item Evaluar experimentalmente la respuesta del sistema bajo diferentes condiciones de funcionamiento, midiendo variables clave como el ángulo de inclinación, la velocidad de respuesta y la efectividad del sistema.
\end{itemize}

\vspace{0.2cm}

Este proyecto puede ser evaluado por un docente de Ciencias Básicas \hspace{0.5cm} SI \underline{\hspace{1cm}} \hspace{0.5cm} NO \underline{\hspace{0.3cm}X\hspace{0.3cm}}

\vspace{0.2cm}

Agradezco la atención prestada.\\
Cordialmente:

\vspace{0.5cm}

\begin{minipage}[t]{0.48\textwidth}
    \rule{6.5cm}{0.4pt}\\
    Nombre: Jeferson Hernán Hernández Moreno\\
    Código: 20201383013\\
    Programa Académico: Ingeniería en Control y Automatización
\end{minipage}
\hfill
\begin{minipage}[t]{0.48\textwidth}
    \rule{6.5cm}{0.4pt}\\
    Nombre: Fabián Andrés Reyes Romero\\
    Código: 20201383011\\
    Programa Académico: Ingeniería en Control y Automatización
\end{minipage}

\vspace{1.5cm}

\begin{center}
    \rule{6.5cm}{0.4pt}\\
    Vo.Bo. Alexander Jiménez Triana, PhD
\end{center}

\end{document}